%==============================================================================
\section{Python Libraries --- Các thư viện Python cho học máy}
%==============================================================================

\begin{dongco}[title={Vì sao dùng thư viện?}]
Thư viện Python là tập các đoạn code và hàm có thể tái sử dụng cho một tác vụ cụ thể.
Chúng giúp giảm công sức lập trình khi tác vụ lặp và phức tạp.
Học máy là lĩnh vực liên ngành: thuật toán kết hợp lập trình và toán.
Thay vì tự cài đặt toàn bộ thuật toán kèm công thức thống kê/toán, ta dùng thư viện để làm nhanh và đúng hơn.
\end{dongco}

\subsection{Danh sách thư viện ML phổ biến}

\begin{phuongphap}[title={Các thư viện được liệt kê}]
\begin{itemize}
  \item NumPy.
  \item Pandas.
  \item SciPy.
  \item Scikit-learn.
  \item PyTorch.
  \item TensorFlow.
  \item Keras.
  \item Matplotlib.
  \item Seaborn.
  \item OpenCV.
  \item NLTK.
  \item spaCy.
  \item XGBoost.
  \item LightGBM.
  \item Gensim.
\end{itemize}
\end{phuongphap}

\subsection{NumPy}

\begin{dinhnghia}[title={NumPy}]
NumPy là package xử lý mảng/ma trận (array/matrix) cho tính toán khoa học:
thao tác đại số tuyến tính, biến đổi Fourier và các phép toán khác.
NumPy cung cấp đối tượng mảng đa chiều hiệu năng cao và công cụ để thao tác ma trận, là thành phần quan trọng của hệ sinh thái ML Python.
\end{dinhnghia}

\begin{phuongphap}[title={Các phép toán quan trọng}]
\begin{itemize}
  \item Phép toán toán học và logic trên array.
  \item Biến đổi Fourier.
  \item Các phép toán liên quan đại số tuyến tính.
\end{itemize}
\end{phuongphap}

\begin{luuy}[title={So sánh trực giác}]
Có thể xem NumPy như “thay thế MATLAB” khi đi cùng SciPy và Matplotlib.
\end{luuy}

\begin{phuongphap}[title={Cài đặt và import}]
\begin{itemize}
  \item Với Anaconda: thường đã có sẵn.
  \item Với Python chuẩn: \texttt{pip3 install numpy}.
  \item Import: \texttt{import numpy as np}.
\end{itemize}
\end{phuongphap}

\begin{vidu}[title={Ví dụ: tạo array 1 chiều}]
\begin{lstlisting}
import numpy as np
data = np.array([1,2,3,4,5])
print(data)
print(len(data))
print(type(data))
print(data.shape)
\end{lstlisting}

\textbf{Output}:
\begin{lstlisting}
[1 2 3 4 5]
5
<class 'numpy.ndarray'>
(5,)
\end{lstlisting}
\end{vidu}

\subsection{Pandas}

\begin{dinhnghia}[title={Pandas}]
Pandas là thư viện mạnh để thao tác và phân tích dữ liệu.
Trong pipeline ML, pandas thường dùng ở bước tiền xử lý/chuẩn bị dữ liệu.
Pandas dựa trên hai cấu trúc dữ liệu chính: \textbf{Series} (1 chiều) và \textbf{DataFrame} (2 chiều).
\end{dinhnghia}

\begin{phuongphap}[title={5 bước xử lý dữ liệu}]
Load, Prepare, Manipulate, Model, Analyze.
\end{phuongphap}

\begin{dinhnghia}[title={Các cấu trúc dữ liệu trong Pandas}]
\begin{itemize}
  \item \textbf{Series}: mảng 1 chiều có nhãn trục (homogeneous).
  \item \textbf{DataFrame}: cấu trúc 2 chiều dạng bảng, có thể chứa dữ liệu khác kiểu (heterogeneous).
  \item \textbf{Panel}: cấu trúc 3 chiều (như container của DataFrame) --- \textbf{đã lỗi thời/deprecated} trong pandas; hiện dùng MultiIndex hoặc xarray.
\end{itemize}
\end{dinhnghia}

\begin{vidu}[title={Ví dụ: Series từ ndarray}]
\begin{lstlisting}
import pandas as pd
import numpy as np
data = np.array(['g','a','u','r','a','v'])
s = pd.Series(data)
print (s)
\end{lstlisting}

\textbf{Output}:
\begin{lstlisting}
0    g
1    a
2    u
3    r
4    a
5    v
dtype: object
\end{lstlisting}
\end{vidu}

\subsection{SciPy}

\begin{dinhnghia}[title={SciPy}]
SciPy là thư viện mã nguồn mở cho tính toán khoa học trên dữ liệu lớn.
Nó bao gồm các module cho tối ưu và các phép toán như tích phân, đại số tuyến tính, xử lý tín hiệu.
SciPy xây trên NumPy nhưng mở rộng để làm các tác vụ phức tạp hơn.
\end{dinhnghia}

\begin{phuongphap}[title={Cài đặt/nhập}]
\begin{itemize}
  \item Với Anaconda: thường đã có sẵn.
  \item Với Python chuẩn: \texttt{pip3 install scipy}.
  \item Ví dụ import: \texttt{from scipy import linalg}.
\end{itemize}
\end{phuongphap}

\begin{vidu}[title={Ví dụ: nghịch đảo ma trận}]
\begin{lstlisting}
import numpy as np
import scipy
from scipy import linalg
A= np.array([[1,2],[3,4]])
print(linalg.inv(A))
\end{lstlisting}

\textbf{Output}:
\begin{lstlisting}
[[-2.   1. ]
 [ 1.5 -0.5]]
\end{lstlisting}
\end{vidu}

\subsection{Scikit-learn}

\begin{dinhnghia}[title={Scikit-learn}]
Scikit-learn là thư viện ML mã nguồn mở xây trên NumPy và SciPy, hỗ trợ supervised và unsupervised learning.
Nó cung cấp công cụ cho tiền xử lý dữ liệu, chọn feature, chọn mô hình, đánh giá mô hình, \ldots
\end{dinhnghia}

\begin{phuongphap}[title={feature nổi bật}]
\begin{itemize}
  \item Xây trên NumPy, SciPy và Matplotlib.
  \item Mã nguồn mở (BSD license).
  \item Dễ tiếp cận và tái sử dụng.
  \item Phủ nhiều thuật toán: classification, clustering, regression, giảm chiều, model selection, \ldots
\end{itemize}
\end{phuongphap}

\begin{phuongphap}[title={Cài đặt/nhập}]
\begin{itemize}
  \item Với Python chuẩn: \texttt{pip3 install scikit-learn}.
  \item Ví dụ import dataset: \texttt{from sklearn.datasets import load\_breast\_cancer}.
\end{itemize}
\end{phuongphap}

\begin{vidu}[title={Ví dụ: load Breast Cancer dataset}]
\begin{lstlisting}
from sklearn.datasets import load_breast_cancer
data = load_breast_cancer()
print(data.target[[10, 50, 85]])
print(list(data.target_names))
\end{lstlisting}

\textbf{Output}:
\begin{lstlisting}
[0 1 0]
['malignant', 'benign']
\end{lstlisting}
\end{vidu}

\subsection{PyTorch}

\begin{dinhnghia}[title={PyTorch}]
PyTorch là thư viện deep learning mã nguồn mở dựa trên Torch, thường dùng phát triển mạng nơ-ron sâu.
Nó dựa trên Python trực quan và có thể định nghĩa computational graph động; phù hợp cho nghiên cứu và phát triển cần linh hoạt.
\end{dinhnghia}

\begin{phuongphap}[title={Cài đặt/nhập}]
\begin{itemize}
  \item Cài (ví dụ CPU Windows, Python 3.8+): \texttt{pip3 install torch torchvision torchaudio}.
  \item Link hướng dẫn cài: \tsurl{https://pytorch.org/get-started/locally/}.
  \item Import: \texttt{import torch}.
\end{itemize}
\end{phuongphap}

\begin{vidu}[title={Ví dụ: NumPy array $\rightarrow$ torch tensor}]
\begin{lstlisting}
import numpy as np
import torch
x = np.ones([3,4])
y = torch.from_numpy(x)
print(y)
\end{lstlisting}

\textbf{Output}:
\begin{lstlisting}
tensor([[1., 1., 1., 1.],
        [1., 1., 1., 1.],
        [1., 1., 1., 1.]], dtype=torch.float64)
\end{lstlisting}
\end{vidu}

\subsection{TensorFlow}

\begin{dinhnghia}[title={TensorFlow}]
TensorFlow là thư viện mã nguồn mở của Google cho ML/DL, hỗ trợ xây computational graph và chạy hiệu quả trên nhiều phần cứng.
Thường dùng cho NLP, nhận dạng ảnh, nhận dạng chữ viết tay.
\end{dinhnghia}

\begin{phuongphap}[title={Cài đặt/nhập}]
\begin{itemize}
  \item Cài: \texttt{pip3 install tensorflow}.
  \item Link hướng dẫn: \tsurl{https://www.tensorflow.org/install/pip}.
  \item Import: \texttt{import tensorflow as tf}.
\end{itemize}
\end{phuongphap}

\begin{vidu}[title={Ví dụ: tạo tensor}]
\begin{lstlisting}
import tensorflow as tf
data = tf.constant([[2,1],[4,6]])
print(data)
\end{lstlisting}

\textbf{Output}:
\begin{lstlisting}
tf.Tensor(
[[2 1]
 [4 6]], shape=(2, 2), dtype=int32)
\end{lstlisting}
\end{vidu}

\subsection{Keras}

\begin{dinhnghia}[title={Keras}]
Keras là thư viện mạng nơ-ron mức cao để tạo mô hình deep learning.
Nó chạy trên TensorFlow/CNTK/Theano và có API trực quan, phù hợp người mới và nghiên cứu vì dễ prototype nhanh.
\end{dinhnghia}

\begin{phuongphap}[title={Cài đặt/nhập}]
\begin{itemize}
  \item Cài: \texttt{pip3 install keras}.
  \item Import: \texttt{import keras}.
\end{itemize}
\end{phuongphap}

\begin{vidu}[title={Ví dụ: dùng CIFAR-10 dataset}]
\begin{lstlisting}
import keras
(x_train, y_train), (x_test, y_test) = keras.datasets.cifar10.load_data()
print(x_train.shape)
print(x_test.shape)
print(y_train.shape)
print(y_test.shape)
\end{lstlisting}

\textbf{Output}:
\begin{lstlisting}
(50000, 32, 32, 3)
(10000, 32, 32, 3)
(50000, 1)
(10000, 1)
\end{lstlisting}
\end{vidu}

\subsection{Matplotlib}

\begin{dinhnghia}[title={Matplotlib}]
Matplotlib là thư viện vẽ phổ biến để trực quan hóa dữ liệu: graph, plot, histogram, bar chart, \ldots
Nó hỗ trợ phân tích, khám phá và trình bày dữ liệu.
\end{dinhnghia}

\begin{phuongphap}[title={Cài đặt/nhập}]
\begin{itemize}
  \item Cài: \texttt{pip3 install matplotlib}.
  \item Import: \texttt{import matplotlib.pyplot as plt}.
\end{itemize}
\end{phuongphap}

\begin{vidu}[title={Ví dụ: vẽ đường thẳng}]
\begin{lstlisting}
import matplotlib.pyplot as plt
plt.plot([1,2,3],[1,2,3])
plt.show()
\end{lstlisting}
\end{vidu}

\subsection{Seaborn}

\begin{dinhnghia}[title={Seaborn}]
Seaborn là thư viện mã nguồn mở xây trên Matplotlib và tích hợp tốt với pandas.
Nó giúp tạo biểu đồ thống kê đẹp và nhiều thông tin, phù hợp phân tích business/marketing.
\end{dinhnghia}

\begin{phuongphap}[title={Cài đặt/nhập}]
\begin{itemize}
  \item Cài: \texttt{pip3 install seaborn}.
  \item Import: \texttt{import seaborn as sns}.
\end{itemize}
\end{phuongphap}

\subsection{OpenCV}

\begin{dinhnghia}[title={OpenCV}]
OpenCV (Open Source Computer Vision Library) là thư viện Python cho thị giác máy tính và xử lý ảnh.
Thư viện này giúp nhận diện pattern trong ảnh, trích xuất feature, và có thể tích hợp với NumPy để xử lý cấu trúc array.
\end{dinhnghia}

\subsection{NLTK}

\begin{dinhnghia}[title={NLTK}]
NLTK (Natural Language ToolKit) là môi trường Python thường dùng phát triển NLP.
Nó có các interface như WordNet và thư viện cho classification, tokenization, parsing, semantic reasoning.
\end{dinhnghia}

\subsection{spaCy}

\begin{dinhnghia}[title={spaCy}]
spaCy là thư viện Python mã nguồn mở, hỗ trợ tác vụ NLP nâng cao theo cách nhanh và hiệu quả.
Tokenization và POS tagging là hai tác vụ thư viện làm tốt.
\end{dinhnghia}

\begin{luuy}[title={Các tool/framework khác}]
XGBoost, LightGBM, Gensim là các công cụ khác trong Python cho ML.
Việc học hệ sinh thái thư viện giúp bạn hiểu cách build/train/deploy mô hình.
\end{luuy}
